% Research proof, 16 October 2026. Two independent internal audits passed.
% No preamble. Label prefix ng-. Uses \norm and standard theorem environments.
\section{Constraint-preserving gauges and coverage of the flat slice}
\label{ng-section}

This module concerns finite compatible periodic regulators, fixed
odd blocking ratio $L$, refinement $M=L^j$, fine spacing $\eta=M^{-1}$,
and final spacing one. Fix $0<\alpha<1$. All constants below are
independent of refinement and volume. We use the corrected nonlinear
map with fixed identity reference, $\mathcal F(A)=\mathcal F^{(j)}(A)$,
and its ordinary potential differential $Q_A=D\mathcal F(A)$.
This is not renamed as the separately based source $Q(U)$.
The flat scalar projection denoted $R$ earlier is denoted $R_0$
here. The potential norms $N_\alpha$, $\mathcal N_\alpha$ and
$\mathcal X_\alpha$ refer to the same complete physical norm,
up to the fixed collar equivalences stated below.

Let $\mathcal X_\alpha$ be the complete-field norm consisting of
the bounded potential, its bounded physical first differences and
the local $\alpha$-H\"older seminorms of those differences. Let
$\mathcal Y_\alpha$ be the scalar $C^\alpha$ norm, including the
supremum, and $\mathcal Z_\alpha$ the scalar $C^{2,\alpha}$ norm,
including all physical differences through order two and the
$\alpha$-H\"older members at order two. Fixed collar enlargements
give equivalent norms. At fixed refinement these are Banach spaces.
The estimates below are uniform in that refinement. For distances
above a fixed collar size, supremum norms supply the corresponding
global increments. No field is sharply cut off before differentiating.

The purpose is to upgrade a local minimum on the fixed flat gauge
slice to a local minimum of the same nonlinear constraint, modulo
constraint-preserving gauge transformations. This requires a new
gauge argument; strict minimality on the slice alone is insufficient.

\subsection{The induced coarse gauge of the corrected average}

For the ordinary group average $M$ and one-step endpoint mean
$a_y(U)$, use the exact identities already checked in Notebook 14:
\[
 C(U)(c)=a_{c_-}(U)^{-1}M(U)(c)a_{c_+}(U),\qquad
 M(U^g)=M(U)^{g|_{\rm coarse}},
\]
where $U^g(e)=g(e_-)U(e)g(e_+)^{-1}$. Direct multiplication gives
\begin{equation}\label{ng-induced-group-gauge}
 C(U^g)=C(U)^{k(U,g)},\qquad
 k_y(U,g)=a_y(U^g)^{-1}g(y)a_y(U).
\end{equation}
All factors are on the common small-field branch. The equality is
exact, not a linearization. In particular $k(U,1)=1$.

To compute its differential, set
\[
 P_x=U(\Gamma_{yx})=e^{iY_x},\quad
 \overline Y_y=\sum_xL^{-d}Y_x,\quad
 a_y(U)=e^{i\overline Y_y},\qquad
 \mathcal J_X Z=\int_0^1\operatorname{Ad}_{e^{itX}}Z\,dt.
\]
For $g(s)=e^{is\lambda}$ the contour left increment is
$\lambda(y)-\operatorname{Ad}_{P_x}\lambda(x)$.
The logarithm and exponential differentials therefore give the
following explicit coarse gauge operator:
\begin{equation}\label{ng-one-step-gauge}
 \begin{split}
 T_U\lambda(y)=\operatorname{Ad}_{a_y(U)^{-1}}
 \bigg[\lambda(y)-\mathcal J_{\overline Y_y}
   \sum_xL^{-d}\mathcal J_{Y_x}^{-1}
       \bigl(\lambda(y)-\operatorname{Ad}_{P_x}\lambda(x)\bigr)
 \bigg].
 \end{split}
\end{equation}
Thus $(1/i)\partial_s k_y(U,e^{is\lambda})|_0=T_U\lambda(y)$.
The gauge parameter has angle units; no factor $h$ is inserted
into $g=e^{i\lambda}$. At $U=1$ this formula reduces exactly to
the normalized scalar block average $Q_{0,h}$.

For $U=e^{ihA}$, every coefficient in
\eqref{ng-one-step-gauge} is a finite local analytic function of
$hA$. Its difference from its value at zero is $O(h\norm{A}_\infty)$.
Finite one-step row incidence and physical-adjoint column incidence
then give
\begin{equation}\label{ng-one-step-smallness}
 \norm{T_U-Q_{0,h}}_{\rm av}\le C_Lh\norm{A}_\infty,
 \qquad \norm{T_U}_{\rm av}\le1+C_Lh\norm{A}_\infty.
\end{equation}
Here the absolute Schur norm uses physical measures and includes
the bounded-field norm of the transpose. Its column factor is the
fixed volume ratio $L^d$. An evaluation at the block center in
\eqref{ng-one-step-gauge} consequently costs a fixed $L$ constant,
not a power of the full refinement. The flat scalar averages and
their physical adjoints have row sums one.

Let $C_t$ be the corrected iterates of $U=e^{i\eta A}$, and let
$g_0=g$, $g_{t+1}=k(C_t,g_t)$. Induction in
\eqref{ng-induced-group-gauge} gives the exact multilevel gauge
law $C^{(j)}(U^g)=C^{(j)}(U)^{g_j}$. Differentiating at $g=1$
gives the scalar operator
\begin{equation}\label{ng-composite-gauge}
 \mathcal G_A=T_{C_{j-1}}\cdots T_{C_0},\qquad
 \mathcal G_0=Q_0.
\end{equation}
These backgrounds are the corrected iterates, unlike the ordinary
reference trajectory used for the source $Q(U)$. The proven bound
$C_t=e^{ih_tA_t}$, $\norm{A_t}_\infty\le2\norm{A}_\infty$,
and $\sum_{t<j}h_t\le1/(L-1)$ imply
\begin{equation}\label{ng-composite-smallness}
 \norm{\mathcal G_A-Q_0}_{\rm av}
       \le e^{C\norm{A}_\infty/(L-1)}-1.
\end{equation}
Its coefficients have fixed final physical range. The same estimate,
up to a fixed constant, holds in weighted physical-cell norms.
The operator is holomorphic in the common bounded-field ball;
local analytic coefficients and finite range first give holomorphy
at each refinement, and the displayed uniform bounds give uniform
Cauchy bounds on smaller balls.

Define the ordinary-coordinate fine gauge generator and its
coarse analogue by
\begin{equation}\label{ng-generators}
 X_A=\mathcal J_{\eta A}^{-1}D_U,
 \qquad U=e^{i\eta A},\quad
 D_U\lambda(e)=\eta^{-1}
       (\operatorname{Ad}_{U_e}\lambda(e_+)-\lambda(e_-)).
\end{equation}
A gauge $e^{is\lambda}$ has ordinary potential velocity
$-X_A\lambda$. At the output, with $B=\mathcal F(A)$ and
$V=e^{iB}$, differentiation of the multilevel group law yields
\begin{equation}\label{ng-gauge-intertwining}
 Q_AX_A=\mathcal J_B^{-1}D_V\mathcal G_A.
\end{equation}
The output spacing is one. In particular $X_A\lambda$ is tangent
to the nonlinear constraint whenever $\mathcal G_A\lambda=0$.

\subsection{Smooth scalar tools for the flat gauge}

Retain $\mathcal L=D^TD$, the scalar flat projection $R_0$, and
$Z$ of the previous flat notebooks. Put
\begin{equation}\label{ng-T-definition}
 T=G'R_0=Z\mathcal L.
\end{equation}
The exact identities give
\begin{equation}\label{ng-T-identities}
 Q_0T=0,\qquad R_0\mathcal L T=R_0.
\end{equation}
The following scalar estimates are needed:
\begin{equation}\label{ng-scalar-bounds}
 \norm{R_0f}_{\mathcal Y_\alpha}\le C\norm{f}_{\mathcal Y_\alpha},
 \qquad
 \norm{Tf}_{\mathcal Z_\alpha}\le C\norm{f}_{\mathcal Y_\alpha}.
\end{equation}
Here is a direct justification from the already proved canceled
aliases. With $A_0=a_0a_0^-$ and
$D_{s,2}=A_0+P_0^2\sum_{r\ne0}a_ra_r^-/P_r^2$, the high block of
$T=Z\mathcal L$ is
\[
 T_{rs}=\frac{\delta_{rs}}{P_r}
       -\frac{P_0^2a_r^-a_s}{P_r^2P_sD_{s,2}},\qquad r,s\ne0.
\]
Its mixed entries are
$T_{0s}=-a_0^-a_s/(P_sD_{s,2})$ and
$T_{r0}=-P_0a_r^-a_0/(P_r^2D_{s,2})$, while
$T_{00}=P_0R_{s,2}/D_{s,2}$.
Subtract $(\mathcal L+1)^{-1}$. The high diagonal remainder has
heat factor $1-e^{-t}$, making its two-output-difference
$\alpha$-H\"older estimate integrable. All high rank output
vectors are $a_r^-/P_r^2$, which are absolutely summable after
$2+\alpha$ differences; input rows are in alias $\ell^2$.
The scalar massive part maps $C^\alpha$ to $C^{2,\alpha}$ by
zero-row heat cancellation and the time split at the square of
the output separation. This is the same second/third heat
derivative calculation used for the right-divergence Schauder
theorem, now with both derivatives on the output.

Also $R_0=I-P_0^{\rm gauge}$, where the rank-one Bloch projection
has output vector
\[
 \chi_0=a_0^-,\qquad \chi_r=P_0a_r^-/P_r\quad(r\ne0),
 \qquad P_0^{\rm gauge}=\chi\chi^-/D_{s,2}.
\]
It maps bounded fine fields to $C^{1,\beta}$ for every fixed
$\beta<1$, by absolute output alias sums and input $\ell^2$.
This proves the first bound in \eqref{ng-scalar-bounds} for fixed
$\alpha<1$. Periodic reconstruction and contour shifting can
be applied to these complete operators; hence the same estimates
have exponential cell locality and weighted versions. These
arguments do not claim a second-difference estimate on arbitrary
bounded forcing for the massive inverse.

We also require a sufficiently smooth scalar right inverse of
$Q_0$. Define $H_b$ by the canceled scalar biharmonic symbols
\begin{equation}\label{ng-biharmonic-extension}
 (H_b)_0=\frac{a_0^-}{D_{s,2}},\qquad
 (H_b)_r=\frac{P_0^2a_r^-}{P_r^2D_{s,2}}\quad(r\ne0).
\end{equation}
Then
\begin{equation}\label{ng-biharmonic-bounds}
 Q_0H_b=I,\qquad
 \norm{H_bc}_{\mathcal Z_\alpha}\le C\norm{c}_\infty.
\end{equation}
On nonzero real fibers this is exactly
$\mathcal L^{-2}Q_0^T(Q_0\mathcal L^{-2}Q_0^T)^{-1}$.
The canceled formula supplies the constant-mode continuation, so
the complete operator has the periodic gluing needed above.
The identity follows by multiplying the scalar row $a_r$ into
\eqref{ng-biharmonic-extension}. The bound follows from the
same $a_r/P_r^2$ absolute alias sums; these even allow $3+\beta$
differences for $\beta<1$. Central entries are bounded and analytic
through $P_0=0$. This is a separately constructed scalar extension,
not the less smooth extension $G'^2Q_0^T(Q_0G'^2Q_0^T)^{-1}$.

\subsection{An invertible gauge-fixing differential}

Set $F_A=\mathcal G_A-Q_0$. On a sufficiently small field ball,
\begin{equation}\label{ng-scalar-chart}
 T_A^s=(I+H_bF_A)^{-1}
\end{equation}
is a uniformly bounded holomorphic operator on $\mathcal Z_\alpha$.
Indeed $F_A:\mathcal Z_\alpha\to\ell^\infty$ has norm
$O(\norm{A}_\infty)$ and $H_b$ satisfies
\eqref{ng-biharmonic-bounds}. The inverse differs from identity
by the same small order. If $Q_0\psi=0$ and
$\lambda=T_A^s\psi$, applying $Q_0$ to
$\lambda+H_bF_A\lambda=\psi$ gives
\begin{equation}\label{ng-scalar-kernel}
 \mathcal G_A\lambda=0.
\end{equation}

There is a uniform estimate
\begin{equation}\label{ng-generator-bound}
 \norm{(X_A-D)\lambda}_{\mathcal X_\alpha}
       \le C\norm{A}_{\mathcal X_\alpha}
                        \norm{\lambda}_{\mathcal Z_\alpha}.
\end{equation}
To see the physical scaling, write
$D_U=D+E_s(A)$ with
$E_s(A)\lambda=\eta^{-1}(\operatorname{Ad}_{e^{i\eta A}}-I)
\lambda(\cdot+\eta e_\mu)$.
The coefficient in this expression extends analytically to
$\eta=0$ and has $C^{1,\alpha}$ norm at most
$C\norm{A}_{\mathcal X_\alpha}$.
Also $\mathcal J_{\eta A}^{-1}-I$ has that norm at most
$C\eta\norm{A}_{\mathcal X_\alpha}$.
Discrete product and shift estimates then prove
\eqref{ng-generator-bound}, using two derivatives of $\lambda$.
Every derivative of a coefficient is paid by a derivative of $A$;
there is no inverse-mesh loss. These maps are holomorphic in their
indicated Banach norms on a common smaller ball.

The operator $R_0D^T:\mathcal X_\alpha\to\mathcal Y_\alpha$
is uniformly bounded. Let
$\mathcal R_\alpha=\{f\in\mathcal Y_\alpha:R_0f=f\}$,
a closed subspace because of \eqref{ng-scalar-bounds}. Define on it
\begin{equation}\label{ng-fixing-map}
 \mathcal M_A=R_0D^TX_AT_A^sT.
\end{equation}
By \eqref{ng-T-identities}, $\mathcal M_0=I$ on
$\mathcal R_\alpha$. The preceding estimates give
\begin{equation}\label{ng-fixing-smallness}
 \norm{\mathcal M_A-I}_{\mathcal R_\alpha\to\mathcal R_\alpha}
       \le C\norm{A}_{\mathcal X_\alpha}.
\end{equation}
Its range is in $\mathcal R_\alpha$ because it begins with $R_0$.
For a fixed small radius its Neumann inverse is uniformly bounded
and holomorphic. Define the scalar correction operator
\begin{equation}\label{ng-correction}
 \mathcal L_A=T_A^sT\mathcal M_A^{-1}R_0D^T.
\end{equation}
It satisfies the exact identities
\begin{equation}\label{ng-correction-identities}
 \mathcal G_A\mathcal L_Av=0,\qquad
 R_0D^TX_A\mathcal L_Av=R_0D^Tv,
\end{equation}
and the stronger norm estimate
\begin{equation}\label{ng-correction-bound}
 \norm{\mathcal L_Av}_{\mathcal Z_\alpha}
       \le C\norm{R_0D^Tv}_{\mathcal Y_\alpha}.
\end{equation}

Consequently any constrained tangent $Q_Av=0$ decomposes as
\begin{equation}\label{ng-tangent-decomposition}
 v=\bigl(v-X_A\mathcal L_Av\bigr)+X_A\mathcal L_Av,
\end{equation}
where the first term is both constrained and in the fixed flat
gauge, and the second is a constraint-preserving gauge tangent.
The first claim uses \eqref{ng-gauge-intertwining} and
\eqref{ng-correction-identities}. Its tangent space is exactly the
image of $X=\ker Q\cap\ker R_0D^T$ under the constraint-chart
derivative $T_A=(I+H D\mathcal R(A))^{-1}$: for a constrained
fixed-gauge tangent $v$, the vector
$z=(I+H D\mathcal R(A))v$ satisfies $Qz=0$ and $R_0D^Tz=0$.
A stationary Wilson field on
the fixed slice is therefore stationary on the full constraint:
the action derivative vanishes on the first term by slice
stationarity and on the second by exact gauge invariance.

\subsection{A constraint-preserving flow into the slice}

For real fields use the differential equation
\begin{equation}\label{ng-flow}
 \frac{dA(t)}{dt}=-X_{A(t)}\mathcal L_{A(t)}A(t),
 \qquad A(0)=A^{\rm in}.
\end{equation}
The vector field is locally Lipschitz on the small
$\mathcal X_\alpha$ ball, uniformly in refinement. Its negative
sign matches the actual gauge velocity in
\eqref{ng-generators}. Along any solution in that ball,
\begin{equation}\label{ng-flow-identities}
 R_0D^TA(t)=e^{-t}R_0D^TA^{\rm in},\qquad
 \mathcal F(A(t))=\mathcal F(A^{\rm in}),\qquad
 \mathcal S_\eta(e^{i\eta A(t)})
             =\mathcal S_\eta(e^{i\eta A^{\rm in}}).
\end{equation}
The first identity follows directly from
\eqref{ng-correction-identities}; the second from
\eqref{ng-gauge-intertwining}; the third from gauge invariance.
All parameter operators preserve real fields. Thus the flow
integrates unitary gauge transformations on the finite regulator.

By \eqref{ng-correction-bound} and the uniform bound for $X_A$,
\begin{equation}\label{ng-flow-speed}
 \norm{A'(t)}_{\mathcal X_\alpha}
    \le Ce^{-t}\norm{R_0D^TA^{\rm in}}_{\mathcal Y_\alpha}.
\end{equation}
Choose the initial ball sufficiently small that its radius plus
this total displacement lies within the established operator
domain. Standard continuation then keeps the solution in that
domain for every $t\ge0$. The integrable speed proves convergence
in $\mathcal X_\alpha$ to $A^\flat$, with
\begin{equation}\label{ng-flow-limit}
 R_0D^TA^\flat=0,\quad
 \mathcal F(A^\flat)=\mathcal F(A^{\rm in}),\quad
 \mathcal S_\eta(e^{i\eta A^\flat})
          =\mathcal S_\eta(e^{i\eta A^{\rm in}}),\quad
 \norm{A^\flat-A^{\rm in}}_{\mathcal X_\alpha}
       \le C\norm{R_0D^TA^{\rm in}}_{\mathcal Y_\alpha}.
\end{equation}
The gauge parameter is integrable even in $\mathcal Z_\alpha$,
so the gauge transformations themselves converge uniformly to a
unitary limit. Hence the limiting field remains in the original
gauge orbit, not merely at the same action value.

For a fixed slice field $A_*$ the same construction can be made
in a neighborhood of $A_*$. Since $R_0D^TA_*=0$,
the total displacement is bounded by
$C\norm{A^{\rm in}-A_*}_{\mathcal X_\alpha}$.
Choosing this neighborhood inside the strict slice-minimality
neighborhood ensures that $A^\flat$ belongs to it. Therefore,
if $A_*$ is a strict local minimum on
\[
 \mathcal F(A)=B,\qquad R_0D^TA=0,
\]
then it is a local minimum of $\mathcal F(A)=B$ without imposing
the flat gauge condition. Equality in a sufficiently small
neighborhood implies $A^\flat=A_*$ and hence that the initial
field is gauge equivalent to $A_*$. This is local minimality
modulo the constraint-preserving gauge orbit, not an absolute
global minimum among all lattice configurations.

The argument supplies uniform local gauge coverage for the stated
global small-field class. It does not cover arbitrary plaquette-small
gauge orbits, changing retained domains, topologically nontrivial
bundles, or a continuum quantum measure. No finite total action is
assigned to arbitrary bounded infinite-lattice fields here.
